\entete{FOD}{12 avril 2012}

\section{Alg\`ebre relationnelle et SQL (10 points)}

Soit la base de donn�es suivante (simplifi�e) qui sert � l'�valuation des projets d�pos�s en r�ponse � un appel d'offres. Chaque projet a un auteur et doit avoir entre 3 et 5 rapporteurs (rapporteur = personne qui pr�pare un rapport d'�valuation). Une personne peut d�poser z�ro, un ou plusieurs projets. Une personne peut �tre rapporteur pour un seul projet. Les conflits d'int�r�t doivent �tre �vit�s : d'abord, une personne ne peut pas �tre rapporteur pour les projets dont elle est l'auteur ; les autres conflits d'int�r�t sont r�pertori�s dans la table CONFLIT : chaque enregistrement dans cette table correspond � une paire [personne1, personne2] telle que personne1 ne peut pas �tre rapporteur pour un projet dont l'auteur est personne2.

\smallskip

\texttt{\textbf{Personne}(id\_personne, nom, pr�nom, date\_naissance)}\\ 
\texttt{\textbf{Projet}(id\_projet, id\_auteur, nom\_projet, description)}\\ 
\texttt{\textbf{Rapport}(id\_rapporteur, id\_projet, texte\_�valuation, date)}\\
\texttt{\textbf{Conflit}(id\_rapporteur, id\_auteur)}\\
\smallskip


\begin{enumerate}

\item Pour chaque table, identifier les cl\'es primaires et les cl\'es \'etrang\`eres. (1 point)
\begin{correction}\\
\begin{itemize}
\item Cl\'es primaires\\
    Table \texttt{Personne : id\_personne}\\
    Table \texttt{Projet : id\_projet}\\
    Table \texttt{Rapport : id\_rapporteur} (vu qu'une personne peut �tre rapporteur pour un seul projet)\\ 
    ou \'eventuellement \texttt{Rapport : (id\_rapporteur, id\_projet)}\\
    Table \texttt{Conflit : (id\_rapporteur, id\_auteur)}\\
\item Cl\'es \'etrang\`eres :\\
  \texttt{Projet[id\_auteur] $\subseteq$ Personne[id\_personne]}\\
  \texttt{Rapport[id\_rapporteur] $\subseteq$ Personne[id\_personne]}\\
  \texttt{Rapport[id\_projet] $\subseteq$ Projet[id\_projet]}\\
  \texttt{Conflit[id\_rapporteur] $\subseteq$ Personne[id\_personne]}\\
  \texttt{Conflit[id\_auteur] $\subseteq$ Personne[id\_personne]}\\
\end{itemize}
\end{correction}

\item Donner des commandes SQL permettant de cr\'eer les tables
\texttt{\textbf{Rapport}} et \texttt{\textbf{Conflit}}. (1 point) 


\begin{correction}
  \begin{lstlisting}
  CREATE TABLE Rapport (
  id_rapporteur INTEGER NOT NULL,
  id_projet INTEGER NOT NULL,
  texte_�valuation VARCHAR2(2000),
  date DATE,
  CONSTRAINT Rapport_PK PRIMARY KEY (id_rapporteur),
  CONSTRAINT Rapport_id_rapporteurKF FOREIGN KEY (id_rapporteur) REFERENCES Personne (id_personne),
  CONSTRAINT Rapport_id_projetKF FOREIGN KEY (id_projet) REFERENCES Projet (id_projet) );
  );

  CREATE TABLE Conflit (
  id_rapporteur INTEGER NOT NULL,
  id_auteur INTEGER NOT NULL,
  CONSTRAINT Conflit_PK PRIMARY KEY (id\_rapporteur, id_auteur),
  CONSTRAINT Conflit_id_rapporteurKF FOREIGN KEY (id_rapporteur) REFERENCES Personne (id_personne),
  CONSTRAINT Conflit_id_auteurKF FOREIGN KEY (id_auteur) REFERENCES Personne (id_personne) );
  \end{lstlisting}

\end{correction}

\item Donner les identifiants des projets pour lesquels au moins un rapport a �t� transmis apr�s le 01-01-2012. Alg\`ebre et SQL (1 point) 
\begin{correction}

  Alg\`ebre:\\ $\pi_{id\_projet}(\sigma_{date \ge 01-01-2012}(Rapport))$\\ 
  SQL: 
   \begin{lstlisting}
    SELECT id_projet FROM Rapport
    WHERE date > 01-01-2012;
    \end{lstlisting}
    
\end{correction}

\item Donner le nom de l'auteur qui a le plus de conflits d'int�r�t. SQL seulement. (2 points) 

\begin{correction}

 \begin{lstlisting}
    SELECT id_personne, nom, count(*)
    FROM Personne, Conflit
    WHERE id_personne = id_rapporteur
    GROUP BY id_personne
    HAVING count(*) >= ALL (SELECT count(*)
                               FROM Conflit
                               GROUP BY id_rapporteur);
 \end{lstlisting}
 
\end{correction}

\item Donner, pour chaque personne (nom et pr\'enom), le nombre de projets dont elle a �t� auteur; les r�sultats doivent �tre tri�s par ordre d�croissant du nombre de projets. SQL seulement (1 point) 
\begin{correction}

 \begin{lstlisting}
  SELECT id_personne, nom, pr�nom, count(id_projet) AS nbProjets
  FROM Personne, Projet
  WHERE id_personne = id_auteur
  GROUP BY id_personne
  ORDER BY nbProjets DESC;
 \end{lstlisting}
 
\end{correction}

\item Donner les identifiants des projets qui ont un nombre de rapports insuffisant (moins de 3). SQL seulement. (2 points) 

\begin{correction}\\ 

  SQL :
   \begin{lstlisting}
   SELECT id_projet, count(id_rapporteur)
   FROM Rapport
   GROUP BY id_projet
   HAVING count(id_rapporteur) < 3;
   \end{lstlisting}

\end{correction}


\item Donner les identifiants des personnes ayant fait un rapport pour chaque projet de Jean DUPONT.
 Alg\`ebre seulement. (1 point) 
\begin{correction}\\
  $R_1$ (les identifiants des projets de Jean DUPONT) :\\
  $\pi_{id\_projet}(Projet \Join_{id\_auteur = id\_personne} (\sigma_{nom='DUPONT' \wedge pr�nom='Jean'} (Personne)))$.
  
  $R_2$ (la correction):\\
  $(\pi_{id\_rapporteur,id\_projet}(Rapport)) \div R_1$
\end{correction}
  
\end{enumerate}


\section{Indexation (3 points)}

On ins\`ere successivement les enregistrements suivants dans la table Projet, selon cet ordre :

\begin{center}
\begin{small}
\begin{tabular}{ccccc}
id\_projet & id\_auteur & \textbf{nom\_projet} & description\\
1 & 5 & AERO & ...\\
2 & 3 & MAXIMIN & ...\\
3 & 10 & NETSET & ...\\
4 & 7 & SYSTEME & ...\\
5 & 4 & STAR & ...\\
6 & 5 & TRECVID & ...\\
7 & 12 & SGAME & ...\\
8 & 4 & INITIAL & ...\\
9 & 2 & METRONOM & ...\\
10 & 1 & VIASAT & ...\\
11 & 3 & HERERA & ...\\
12 & 5 & DOCEND & ...\\
13 & 8 & DIMINUTIS & ...\\
14 & 9 & EXSAT & ...\\
15 & 10 & DATASEC & ...\\
16 & 2 & LEMME & ...\\
17 & 12 & VISTARA & ...
\end{tabular}

\end{small}
\end{center}

On met en place un index sur l'attribut \textbf{nom} de cette table. 

\begin{enumerate}
\item Construire l'arbre B+ d'ordre \textbf{2} correspondant \`a cet ensemble d'articles, en respectant l'ordre d'arriv\'ee. Donner les \'etapes de construction interm\'ediaires importantes (celles qui produisent un changement de la structure de l'arbre) et distinguer sur les sch\'emas les n{\oe}uds de natures diff\'erentes (n{\oe}uds contenant des cl\'es vs. ceux contenant les enregistrements complets ou autre chose). (2 points)

\begin{correction}

La figure \ref{fi:sol_Bp} donne l'arbre B+ d'ordre 2 obtenu. Ici les n{\oe}uds internes (en italique) contiennent les cl\'es des enregistrements, alors que les n{\oe}uds feuilles (en gras) contiennent les enregistrements complets, ou bien leurs cl\'es plus les pointeurs vers ces enregistrements sur disque.

\begin{figure}[htbp]
\begin{center}
\epsfig{file=sol_Bp.eps, width=15cm}
\end{center}
\caption{Arbre B+ d'ordre 2.}
\label{fi:sol_Bp}
\end{figure}

\end{correction}

\item On recherche les auteurs des projets dont l'initiale du nom se trouve entre 'D' et 'I' (inclus). Combien de pages disque doivent \^etre charg\'ees pour r\'epondre \`a cette requ\^ete en utilisant l'arbre B+ ? Justifier. (1 point)

\begin{correction}

Il s'agit ici d'une requ\^ete par intervalle, facilement r\'ealisable avec un arbre B+. On commence par faire une recherche du premier enregistrement dont l'initiale du nom est \'egale ou imm\'ediatement sup\'erieure \`a 'H'. On charge le n{\oe}ud racine 'Nelson' puis son fils gauche puisque 'H' pr\'ec\`ede 'N'. 'H' se situant entre 'Eames' et 'Jacobsen', on lit ensuite le n{\oe}ud feuille relatif aux enregistrements 'Gray' et 'Jacobsen'. Jusque l\`a trois n{\oe}uds, et donc trois pages, ont \'et\'e charg\'es. Dans cette derni\`ere page charg\'ee, 'Jacobsen' est s\'electionn\'e car plus grand que 'H' tout en \'etant plus petit que 'N'. On a trouv\'e le premier enregistrement r\'epondant \`a la requ\^ete, il suffit d'exploiter le cha\^inage des feuilles pour trouver les autres enregistrements : on lit ainsi les deux n{\oe}uds feuilles suivants, jusqu'\`a l'enregistrement 'Panton' qui ne rentre par dans l'intervalle recherch\'e. On a trouv\'e cinq enregistrements r\'epondant \`a la requ\^ete ('Jacobsen', 'Jouin', 'Le Corbusier', 'Nelson' et 'Noguchi'), et lu deux pages disque de n{\oe}uds internes, plus trois pages disque de n{\oe}uds feuilles. Si les feuilles contiennent les enregistrements complets, alors les pr\'enoms recherch\'es y figurent et on a fini : {\bf on a lu cinq pages disque en tout}. Sinon, il faut acc\'eder au fichier des donn\'ees \`a partir des pointeurs contenus dans les feuilles. Dans ce cas on doit lire au minimum une page disque (cas optimiste o\`u les cinq enregistrements s\'electionn\'es sont tous sur la m\^eme page) et au maximum cinq pages (cas pessismiste o\`u ils sont tous r\'epartis sur des pages diff\'erentes) : {\bf on aura donc lu entre six et dix pages en tout}.

\end{correction}

\end{enumerate}


\section{Optimisation (3 points)}

\begin{enumerate}
 \item Donner le plan EXPLAIN obtenu pour la requ\^ete suivante~(2 points)~:\\
\hspace*{1cm}SELECT pe.nom, pe.prenom\\
\hspace*{1cm}FROM Personne pe, Projet pr\\
\hspace*{1cm}WHERE pr.id\_auteur = pe.id\_personne AND\\
\hspace*{2cm}pr.nom\_projet = 'METRONOM' ;\\
Sachant que nous avons qu'un index BTREE\_NOM\_PROJET sur le nom du projet (comme cr\'e\'e dans la partie~2), le nom n'est pas forc�ment unique.
Vous ferez appara\^itre les conditions de la requ\^etes sous le plan EXPLAIN (comme indiqu\'e dans le cours). Expliquer votre choix.

\begin{correction}
\begin{verbatim}
0  SELECT STATEMENT
1    NESTED LOOP JOIN
2      TABLE ACCESS BY ROWID    Projet           pr
3*       INDEX RANGE SCAN       BTREE_NOM_PROJET
4*     TABLE ACCESS FULL        Personne         pe

(3) nom_projet = METRONOM
(4) id_auteur = id_personne
\end{verbatim}
\end{correction}

 \item Sachant que nous ajoutons un index de type B+Tree sur la cl\'e primaire \textbf{Personne.id\_personne}, modifier le plan EXPLAIN pr\'ecedent pour y int\'egrer cet index. Expliquer votre choix (1 point)\,:

\begin{correction}
\begin{verbatim}
0  SELECT STATEMENT
1    NESTED LOOP JOIN
2      TABLE ACCESS BY ROWID    Projet           pr
3*       INDEX UNIQUE SCAN      BTREE_NOM_PROJET
4      TABLE ACCESS BY ROWID    Personne         pe
5*       INDEX UNIQUE SCAN      BTREE_ID_PERSONNE

(3) nom_projet = METRONOM
(5) id_auteur = id_personne
\end{verbatim}
\end{correction}

\end{enumerate}

\section{Concurrence (4 points)}

\begin{enumerate}
  \item Indiquez si les ex\'ecutions concurrentes suivantes peuvent produire
des anomalies (et lesquelles), voire \^etre s\'erialisable. (1~point)

\begin{itemize}
  \item $r_1[x]w_1[x]r_2[x] w_2[y] R_1 C_2$
  \item $r_1[x]w_1[x]r_2[y] w_2[y] R_1 C_2$
 \item $r_1[x]r_2[x]r_2[y] w_2[y]r_1[z] R_1 C_2$
  \item $r_1[x]r_2[x]w_2[x] w_1[x] C_1 C_2$
  \item $r_1[x]r_2[x]w_2[x] r_1[y] C_1 C_2$
 \item $r_1[x]w_1[x]r_2[x] w_2[x] C_1 C_2$
\end{itemize}

\begin{correction}

\begin{itemize}
  \item Lecture sale de $T_2$ sur $T_1$.
  \item Aucun conflit.
 \item Aucun conflit.
  \item Risque de mise � jour perdue.
  \item Ex\'ecution s\'erialisable.
 \item Idem.
\end{itemize}
\end{correction}

\item Indiquez si l'ex\'ecution suivante est s\'erialisable. Justifiez votre r\'eponse en donnant les conflits et un graphe de s\'erialisation. (1 point)
\begin{itemize}
  \item
  $r_1[x]r_2[x]w_2[x] C_2 r_3[x]r_4[z]w_1[x]w_3[y]w_3[x]
  C_3 w_1[y]w_5[x]w_1[z] C_1 w_5[y]r_5[z] C_5 C_4$
\end{itemize}

\begin{correction}
  \begin{itemize}
    \item Conflits cycliques entre $T_2$ et $T_3$.
  \end{itemize}
\end{correction}

  \item Faire l'ex\'ecution pr\'ec\'edente avec  un protocole de
  verrouillage � deux phases. Donner le d�roulement de l'histoire qui en r�sulte (avec les verrous), et les interblocages s'ils existent (\textit{deadlock})\,? (2 points)

\end{enumerate}
